\documentclass[11pt]{article}
\usepackage[a4paper,margin=25mm]{geometry}
\usepackage{amsmath,amssymb,amsthm,booktabs}
\newtheorem{theorem}{Theorem}\newtheorem{lemma}[theorem]{Lemma}
\newcommand{\Aut}{\operatorname{Aut}}\newcommand{\Cor}{\operatorname{Cor}}
\title{Coset coordinates and the last six Higman--Sims geometry groups}
\author{}\date{}
\begin{document}\maketitle
Let $G=\mathrm{HS}$, with $|G|=44\,352\,000$, and let $A=\Aut(G)$,
with $|A|=88\,704\,000$. We acquire the reviewed faithful models on
100 points, including $G'=A'=G$ and $C_A(G)=1$. The fixed source is the
23-page supplement dated 6 October 2003 accompanying D.~Leemans,
\emph{The residually weakly primitive geometries of HS}, Australasian
Journal of Combinatorics 33 (2005), 231--236. Its classification and the
exact correspondence between source diagrams and saved parabolic families
are premises of this application.

\begin{theorem}
For the six specified families, the full groups are as follows. Here $n$
is the smallest type size, $a$ lists every model type of that size, $d$
is the valency of the intrinsic graph, and $q$ is the order of its full
automorphism group divided by $|G|$.
\begin{center}\small
\begin{tabular}{cccrrcll}\toprule
Family&Source&$a$&$n$&$d$&$q$&$\Aut(\Gamma)$&$\Cor(\Gamma)$\\\midrule
1&2.6&2&36960&24&2&$A$&$A$\\
2&2.5&1&11550&64&2&$A$&$A$\\
10&3.27&2,3&15400&180&1&$G$&$A$\\
16&3.21&2,3&15400&24&2&$A$&$A\times C_2$\\
17&3.20&1,3&15400&24&2&$A$&$A\times C_2$\\
18&3.19&1&15400&24&2&$A$&$A$\\\bottomrule
\end{tabular}
\end{center}
In the retained source, 3.20 prints $\Aut(\Gamma)=G$ and
$\Cor(\Gamma)=A$, while 3.21 prints $\Aut(\Gamma)=G$ and
$\Cor(\Gamma)=G\times C_2$. The theorem replaces both entries in each
of these two diagrams. Pages 7 and 8 were inspected directly. The four
other rows agree with the printed groups, using the source convention
when a correlation group is omitted. Previously recorded corrections
to a type size in 2.6 and an incidence number in 2.5 are separate.
\end{theorem}

\paragraph{Proving the point action from literal cosets.}
Fix a family $P_1,\ldots,P_r$ and a smallest type $a$. Put $P=P_a$ and
$n=[G:P]$. Let $g_1,\ldots,g_t$ be the ordered generators of the original
100-point model. The certificate supplies elements $c_1,\ldots,c_n\in G$
and permutations $r_1,\ldots,r_t$ of $\{1,\ldots,n\}$. The reader verifies
\begin{equation}\label{edges}
 c_i g_j c_{i^{r_j}}^{-1}\in P\qquad(1\le i\le n,\ 1\le j\le t).
\end{equation}
It also verifies transitivity of $\langle r_1,\ldots,r_t\rangle$, the
existence of a label $b$ with $c_b\in P$, and $\operatorname{Core}_G(P)=1$
in the original model.

The map $i\mapsto Pc_i$ respects each generator by~\eqref{edges}.
Every coset is reached from $P$ by a positive generator word, since $G$
is finite. Thus this map is surjective; its domain and codomain both have
$n$ elements, so it is bijective. Under this bijection the $r_j$ are
exactly right multiplication by $g_j$. Consequently they define a
homomorphism $\phi:G\to E=\langle r_1,\ldots,r_t\rangle$, whose kernel
is $\operatorname{Core}_G(P)=1$. The image is $E$ by definition.
Only after these checks does the program install the established
homomorphism and its injectivity, surjectivity and order attributes in
GAP. There is no unchecked homomorphism premise. Across the nine anchors,
all $2\,813\,580$ generator edges are checked.

\newpage
\paragraph{Recovering objects from points.}
Write $E_b=\phi(P_a)$. A type-$j$ base object has point shadow
$S_j=b^{\phi(P_j)}$. Direct shadows suffice except at the two anchors
of family 10, where type 1 requires additional incident objects.
An encoding consists of its shadow together with typed sets of previously
encoded incident objects. The reconstruction order is acyclic. For a
reference to type $k$, the reader constructs the entire orbit of its base
encoding under $\phi(P_j)$ and checks that its size is
$[P_j:P_j\cap P_k]$. Induction then identifies that orbit with all
incident type-$k$ objects, with none omitted or duplicated.

\begin{lemma}\label{root}
Let a group $E$ act on a set of objects equipped with an equivariant
carrier in a point set. Suppose $b$ belongs to the carrier $S$ of an
object $e$, and $K\le E$ fixes $e$ and is transitive on $S$. Then
\[
 E_e=\langle K,(E_b)_e\rangle.
\]
\end{lemma}
\begin{proof}
If $h\in E_e$, then $b^h\in S$. Choose $k\in K$ with $b^k=b^h$.
The element $hk^{-1}$ fixes both $b$ and $e$, so $h\in(E_b)_eK$.
The reverse inclusion follows because both factors fix $e$.
\end{proof}

For each recorded encoding the program checks its exact shadow, every
referenced orbit, invariance under all generators of $\phi(P_j)$, and
membership in $\phi(P_j)$ of every generator of its stabilizer inside
$E_b$. Lemma~\ref{root} therefore gives full stabilizer $\phi(P_j)$.
The orbit of the encoding has size $[G:P_j]$, so the map from type-$j$
objects to encodings is injective. At family 10, anchor 2 reconstructs
type 1 using types 2 and 3, and anchor 3 reconstructs type 1 using type 2.
All remaining encodings use their direct shadows. The verifier requires
every type to be reached. A typed automorphism fixing all anchor points
then fixes every encoding, by induction, and hence every object. Its
restriction to the anchor is faithful.

\paragraph{The intrinsic relation.}
We acquire the previously proved common-incidence formula. If $O$ is a
non-diagonal orbit of $E_b$ on points, the number of type-$j$ objects
incident with a specified pair in the corresponding orbital is
\[
 m_j(O)=\frac{[G:P_j]|S_j|\,|S_j\cap O|}{n|O|}.
\]
This formula counts objects, including distinct objects with equal shadows.
The proof counts triples of an object and an ordered pair in its shadow;
transitivity of $P_j$ on that shadow makes the number of second points
constant. For each pair take the sorted multiset of $([G:P_j],m_j(O))$
over all types. Correlations preserve it, including multiplicities. Order
the distinct off-diagonal signatures lexicographically and choose the
first whose relation has the indicated valency $d$. Every chosen relation
is checked to be symmetric and irreflexive. The selection rule is intrinsic,
so a correlation taking one anchor to another takes its graph to the
other anchor graph. All smallest types are included.

The base neighbour set is propagated along the verified point generators.
The program visits all $n$ points and checks every adjacency row under
every generator afterwards. Thus possible conflicting propagation paths
cannot silently produce another graph. All rows have valency $d$.

\newpage
\paragraph{Complete bounds and an explicit outer action.}
Let $D$ be the full automorphism group of one of these graphs. A complete
unseeded Traces calculation supplies $|D|=q|G|$ as in the table. The graph
reader checks the whole adjacency array and all known $E$ generators
before calling Traces; it checks every returned generator afterwards.
Only exact integer orders below $2^{53}$, with zero exponent and clean
completion status, are accepted. This is a computational upper bound,
not a deduction from the returned generators alone.

When $q=1$, $E\le D$ and equal orders imply $D=E$. When $q=2$, we
identify the extension directly. In the original model put
$B=\bigcap_jP_j$ and $S=\bigcap_jN_A(P_j)$. The calculation verifies
$S\cap G=B$, $|S|=2|B|$, and an element $x\in S\setminus G$ with
$x^2\in P_a$. The latter element exists in each of the seven instances
where $q=2$. Starting with $b^u=b$, it propagates a point permutation $u$
using the equations
\[
 (i^{\phi(g_j)})^u=(i^u)^{\phi(g_j^x)}.
\]
Every encountered equation is checked, every point is reached, and the
result is required to be bijective. The verifier then checks, as literal
permutation identities,
\[
 \phi(g_j)^u=\phi(g_j^x)\quad\hbox{for every }j,
 \qquad u^2=\phi(x^2).
\]
It also checks $u$ on every adjacency row of the actual graph.

These identities show that $u$ normalizes $E$ and $u^2\in E$. Moreover,
$u\notin E$. Otherwise $u=\phi(g)$ would give the same automorphism of
$G$ as $x$, so $xg^{-1}\in C_A(G)=1$, contrary to $x\notin G$.
Consequently $F=\langle E,u\rangle$ has order $2|G|$. Since $F\le D$
and the complete upper bound gives that same order, $D=F$. In both cases
$E\lhd D$ and $D/E$ is abelian. Perfection of $G$ yields
$E=E'\le D'\le E$, hence $D'=E$.

\paragraph{From graph groups to geometry groups.}
The preceding equality identifies the specified $G$ intrinsically at
every possible image anchor. If $c$ is a correlation, conjugation by its
point restriction carries one graph's derived group to the other.
For every $g\in G$, the restriction of $cgc^{-1}$ to the image anchor
therefore equals that of some $h\in G$. The typed automorphism
$h^{-1}cgc^{-1}$ fixes that anchor pointwise and is trivial by the faithful
restriction just proved. Thus $G$ is normal in the full correlation group.

The acquired coset normalizer theorem is now applicable. Let $T$ be the
setwise stabilizer of the parabolic family in $A$, allowing all type
permutations. It gives
\[
 |\Aut(\Gamma)|=[G:B]|S|,\qquad
 |\Cor(\Gamma)|=[G:B]|T|.
\]
For each type permutation, successive exact transporter intersections
retain precisely the automorphisms carrying all already assigned
parabolics to their targets. Positive representatives are checked on the
entire family. Necessary subgroup-order tests discard only impossible
permutations; no successful representative is inferred from an order.

\newpage
\begin{center}
\begin{tabular}{crrrrr}\toprule
Family&$|B|$&$|S|$&$|T|$&$|T\cap G|/|B|$&$|\Cor(\Gamma)|$\\\midrule
1&200&400&400&1&88704000\\
2&240&480&480&1&88704000\\
10&16&16&32&1&88704000\\
16&32&64&128&2&177408000\\
17&32&64&128&2&177408000\\
18&40&80&80&1&88704000\\\bottomrule
\end{tabular}
\end{center}
In each row $S\cap G=B$. The typed group therefore acts faithfully by
conjugation on $G$ and is identified with $G$ or $A$ by its order and
its actual embedding in $\Aut(G)$. The full centralizer is the acquired
quotient $(T\cap G)/B$. In families 16 and 17 it has order two, is central,
and meets the typed group trivially. Their product has the full correlation
group order, proving $A\times C_2$. In the other four rows that centralizer
is trivial; conjugation faithfully embeds the full group in $A$ and its
order identifies it with $A$. These are identifications by explicit
normalizer and centralizer arguments, not order-based group recognition.

\paragraph{Controls and failed computations.}
Literal subgroup enumeration in $S_3,D_8,A_4$ checks all coset edges, the
carrier stabilizer lemma, and the incidence formula against direct counts,
including repeated shadows. A separate example shows that carrier
transitivity cannot be omitted from Lemma~\ref{root}. Actual damaged
certificates change a coset representative, a point generator, a shadow,
and an acyclic type reference; each must be rejected. Graph controls
change domains, edges and permutations. Ordinary and optimized Python
must agree on all mathematical controls.

The Traces interface is additionally checked on every graph on four
labelled vertices against all 24 permutations, including equality with
the group generated by its output. Nonempty input generator lists are
forbidden. An earlier interface that supplied known generators failed
this control: for the complete graph on four vertices with one supplied
transposition it reported order 18 instead of 24. That program, its input
and the diagnostic are retained and supply no upper-bound premise here.
All nine accepted bounds are fresh unseeded computations. Earlier capped
nauty attempts, an overrestrictive preliminary outer-normalizer test and
an invalid time-cap request are also retained without mathematical credit.

\paragraph{Attribution and remaining scope.}
The original geometry classification and source assertions are attributed
to the retained manuscript and supplement. The coordinate correspondence,
shadow reconstruction framework, common-incidence count and coset
normalizer theorem are acquired from separately reviewed work. This
component contributes the literal coset-edge proof, the use of a point
stabilizer to test an object's full stabilizer, the explicit outer action,
and their nine concrete applications. These are elementary arguments;
historical novelty is not claimed. GAP's exact group algorithms and
Traces's complete upper bounds remain computational premises. No Lean
verification of this component is claimed.

The six full-group determinations complete this particular part of the
86-family analysis, conditional on the acquired classification. Abstract
parabolic structures, the other uncompleted source clauses, and final
whole-manuscript review remain separate. There is no new conjecture
resolution or compression measurement in this component.
\end{document}
